\paragraph{IV 16, definitions 1--4 and the intended placement (cote 25, sites 139--141).}
The packet is headed ``Ensemble singulier et supersingulier d'une fonction, et critères différentiels'' and expressly says that the new number will be used in the paragraph on hyperplane sections, while its natural place is paragraph 16. Site 139 is an exceptionally faint carbon impression: its numerous doubtful readings remain independent of the clearer page that follows it. On site 140, definition 1 starts from a regular prescheme, a function $\phi$, and the class of $\phi_x$ in $\mathfrak m_x^2/\mathfrak m_x^3$. It defines the singular and supersingular zeros by the order-two condition and the degeneracy of the resulting quadratic form; the word ``non'' before ``dégénérée'' is deleted. The same page calls a singular zero that is not supersingular an ordinary singular zero. Remark 2 discusses the relation with regularity of $V(\phi)$, then a substantial passage involving multiplicity is crossed out. Definition 3 introduces the geometric version over a field. Its actual typing has $k'$ in ``extension $k'$ de $k$'', but the following line has $k$ followed by a low handwritten comma, not a raised prime; the reading preserves both. Remarks 4 on site 141 discuss the field-valued test, the relative version fiber by fiber, and invariance under multiplication of $\phi$ by a unit, which motivates extending the definitions to a section of an invertible Module. Several base letters there remain overwritten or uncertain between $S$ and $Y$. These definitions, corrections and uncertainties can be checked at \href{https://grothendieck.umontpellier.fr/web/viewer.html?file=../25.pdf#page=140}{site 139}, \href{https://grothendieck.umontpellier.fr/web/viewer.html?file=../25.pdf#page=141}{site 140} and \href{https://grothendieck.umontpellier.fr/web/viewer.html?file=../25.pdf#page=142}{site 141}.

\paragraph{IV 16, principal parts, discriminant and propositions 5--7 (sites 141--148).}
The construction proceeds through the sections $d^i\phi$ of the sheaves of principal parts. The published reference for principal parts is \href{https://www.numdam.org/articles/10.1007/BF02732123/}{EGA IV 16.3.1--16.3.2}; the use of the diagonal there gives the context for these differential constructions. Proposition 5 identifies the zero and singular-zero loci using orders zero and one; its proposed order-two clause is deleted. The surrounding rectangle on site 141 carries the instruction to keep the adjacent text, so the rectangle itself does not delete its contents. Site 142 defines the closed singular subscheme from the first differential and obtains a quadratic section $M(\phi)$ on it from the second principal part. The manuscript introduces the corresponding map from the tangent Module to the differential Module and the determinant $D(\phi)$; site 143 then designates the closed subscheme cut out by this determinant as the supersingular subscheme and leaves a bare heading for proposition 6, asking that the construction be collected there. The next argument assumes the stated smoothness and relative-dimension hypotheses on $X/Y/S$ and on $V^1=V(\phi)^{\mathrm{sing}}$. Its diagram (D) is built from the transitivity sequence and the conormal/differential sequence; site 144 explicitly refers the latter to paragraph 16 and invokes the Jacobian criterion. Proposition 7 on site 145 states the description of the supersingular locus by ramification of $V^1\longrightarrow Y$. The following argument seeks the corresponding identity of closed subschemes. It compares the ideals attached to two maps in a diagram with locally direct summands $P,Q\subset M$, ranks satisfying $p+q=m+1$, and the determinant line $L=\det P\otimes\det Q\otimes\det M^{-1}$. Site 146 supplies the square with its single $L$ object and the two maps $a,b$; the text explicitly leaves a possible sign dependent on conventions. Site 147 identifies the relevant component of $a$ with the earlier discriminant, again with that sign reservation, and asks for a recapitulative proposition for a section of an invertible Module. Site 148 supplies the wholly handwritten diagram (D') for that extension. The author also discusses isolation of ordinary singularities, refers to the Cohen--Macaulay property and excellence, and asks where the isolation result should be put in Chapter IV. The argument's written coefficients and indices are retained: site 144 has $\mathcal O_X$ in the first displayed source, says that $W$ is smooth over $X$, and repeats $\Omega^1_{X/S}$ at the end of the subsequent conormal display; its principal-parts decomposition and quadratic expression use $X/S$ where diagram (D) uses $X/Y$. Site 145 changes rank letters between the exterior-power expressions and has a visible $M$ in the displayed Hom target. Site 146 leaves the connecting signs blank in the sample exact sequence and determinant identity. Site 147 lacks the dual sign on the later intersection's $P^1$ and writes $\Omega^1_{X/Y}+\Omega^1_{Y/S}$ in its second exact sequence. Site 148 prints the dimension thresholds in terms of $\dim_x X$ and immediately gives the plane example $\phi=xy$; the edition preserves this juxtaposition. The tensor exponent on $K$ in diagram (D') remains doubtful, and the horizontal endpoints are solid dots while the vertical endpoints are open zeros. These are observations of the written argument, not completed corrections. See \href{https://grothendieck.umontpellier.fr/web/viewer.html?file=../25.pdf#page=143}{sites 142--143}, \href{https://grothendieck.umontpellier.fr/web/viewer.html?file=../25.pdf#page=145}{sites 144--146} and \href{https://grothendieck.umontpellier.fr/web/viewer.html?file=../25.pdf#page=148}{sites 147--148}.

\paragraph{IV 16, augmented rings and the passage to several equations (sites 149--152).}
The author now explicitly says that one equation is still too particular and proposes an homomorphism $E\longrightarrow\mathcal O_X$ with $E$ locally free of rank $m$. Site 149 introduces a regular augmentation $A\longrightarrow A/I$ by projectivity of $I/I^2$ and the symmetric-algebra comparison with $\mathrm{gr}_I(A)$, while asking whether the terminology adopted in paragraph 19 requires ``quasi-régulière''. The terminology can be compared with \href{https://www.numdam.org/articles/10.1007/BF02732123/}{published EGA IV 16.9.1--16.9.3}, which defines regular and quasi-regular ideals and immersions and uses the symmetric-algebra comparison in 16.9.3. It then defines an augmentation of ordinary quadratic type through a nondegenerate invertible relation Module in degree two. Proposition 8 and lemma 8.1 on site 150 formulate the graded projectivity and rank criteria and discuss multiplication by the quadratic element; the author observes that the lemma's proof uses nonvanishing in the reduced fibers and suggests a more general formulation. The edition preserves both occurrences of the displayed binomial difference with lower index $d$, while the rank parameter is also called $d$; it also preserves the later wording that calls $B$ ``le spectre d'un corps'' after introducing it as a ring. Proposition 9, beginning on site 150 and continuing on site 151, treats the quotient by the image of a projective Module and its induced augmentation. It expresses the condition through the kernel of the first-order map and the induced quadratic map, then writes out how this latter map is obtained from multiplication in the associated graded algebra. One formula runs beyond the right image edge and is not completed editorially. Corollary 10 passes to a regular local ring and gives the formulations in terms of successive equations, the first-order rank and a nondegenerate quadratic element. Its continuation on site 152 explicitly refers to definition 1. The text then introduces singular, ordinary singular and supersingular zeros for the homomorphism $u$, using $\mathcal O_X$-regularity and the local ring of $V(u)$, and mentions the geometric versions over a field. The underlined capital $M$ at the target of $u$ in corollary 10 is retained as written. The long diagonal strokes across this packet have an unresolved cancellation or working-mark function: site 150's stroke stops before proposition 9, site 151 has two separate strokes, and site 152's stroke stops above the later terminology discussion. The text remains present with visible uncertain marginal marks, and the local cancellations remain separately encoded. See \href{https://grothendieck.umontpellier.fr/web/viewer.html?file=../25.pdf#page=150}{site 149}, \href{https://grothendieck.umontpellier.fr/web/viewer.html?file=../25.pdf#page=151}{site 150}, \href{https://grothendieck.umontpellier.fr/web/viewer.html?file=../25.pdf#page=152}{site 151} and \href{https://grothendieck.umontpellier.fr/web/viewer.html?file=../25.pdf#page=153}{site 152}.

\paragraph{IV 16, the differential construction for a Module and its unfinished continuation (sites 153--156).}
Site 153 explains the dependence on the local ring of $V(u)$ apart from the dimension/regularity condition, discusses the rational-point and relative versions, and notes the equivalent presentation of a section through the dual Module. It constructs $V^1\subset V^o=V(u)$ from the maximal exterior power of the first differential. Site 154 explicitly applies proposition 9, reduces the ordinary-singularity test to order two, and introduces the further rank-drop locus $V''$, the kernel line $K$ and cokernel $W$ on $V^1-V''$, the quadratic map $K\longrightarrow\mathrm{Sym}^2(W)$, and its discriminant $D$. Site 155 puts $V^2=V(D)$, records the ordinary-singular locus as $V^1-V''-V^2$, and interposes a local-freeness lemma to justify the construction. Its other displayed expression, for the supersingular locus, contains an actual blank between $V''$ and $V^2$; no union or difference sign has been supplied. The lemma begins with the letter $v$ and subsequently uses $u$. The surrounding pages repeatedly replace typed $S$ by handwritten $Y$, while other occurrences of $S$ remain untouched; site 154 changes rank letters from $n$ to $d$ but leaves $n$ in the discriminant exponent, whose later sign remains overwritten and uncertain. These incomplete revisions are preserved. The final discussion asks for generalizations of propositions 6 and 7, stipulates the smoothness and dimension data for $V^o$ and $V^1$, and explains the intended minimum-dimension choice. On site 156 the proposed comparison with the ramification subscheme is expressly left to be verified: the author says he has not carried out the exercise in detail and points back to the dévissage after proposition 7. The final paragraph records a rank-locus codimension assertion, while its proposed induction explanation is canceled. The guide supplies neither that missing verification nor a new theorem. See \href{https://grothendieck.umontpellier.fr/web/viewer.html?file=../25.pdf#page=154}{site 153}, \href{https://grothendieck.umontpellier.fr/web/viewer.html?file=../25.pdf#page=155}{site 154}, \href{https://grothendieck.umontpellier.fr/web/viewer.html?file=../25.pdf#page=156}{site 155} and \href{https://grothendieck.umontpellier.fr/web/viewer.html?file=../25.pdf#page=157}{site 156}.
